% ============================================================================
%   CHAPTER 1 — INTRODUCTION
%   Aivancity: 3-5 pages, funnel structure (general → specific)
% ============================================================================

\chapter{Introduction}
\label{chap:introduction}

% ----------------------------------------------------------------------------
\section{Macro Context: AI in Healthcare}
\label{sec:intro-macro}

Artificial intelligence has moved from a peripheral tool to a central 
component of many medical workflows over the past decade. The 
transformation has been particularly rapid since 2022, driven by the 
release of general-purpose large language models capable of processing 
free-text clinical documentation and by the emergence of medically 
specialized variants trained on biomedical corpora. Applications now 
range from automated summarization of patient records to differential 
diagnosis, from radiology report generation to the identification of 
adverse drug interactions.

This expansion carries both promise and risk. On the one hand, LLM-based 
tools address longstanding bottlenecks in clinical practice, where the 
volume of unstructured text and the complexity of medical reasoning 
have historically limited the reach of algorithmic support. On the 
other hand, the confidence that clinicians and patients place in these 
tools depends critically on the reliability of the underlying 
probabilistic outputs. A model that is accurate on average but 
overconfident on the cases where it is wrong is worse than a model that 
is less accurate but honestly uncertain: the former induces misplaced 
trust, while the latter enables appropriate deference to human 
judgment. This asymmetry between accuracy and reliability motivates the 
present thesis.

% ----------------------------------------------------------------------------
\section{Micro Context: Clinical Trial Matching}
\label{sec:intro-micro}

Among the many applications of AI in healthcare, patient matching to 
clinical trials occupies a particularly consequential position. Clinical 
trials are the mechanism through which new therapies are evaluated and 
made available to patients, and their success depends on the timely 
recruitment of suitable participants. Despite the substantial resources 
invested in recruitment, less than $5\%$ of adult cancer patients enroll 
in a clinical trial, and $80\%$ of trials fail to meet their recruitment 
targets~\cite{jin2024trialgpt}. The bottleneck lies partly in the sheer 
volume of open studies --- over $400{,}000$ trials are currently 
registered on ClinicalTrials.gov --- and partly in the complexity of 
their eligibility criteria, which combine numeric thresholds, medical 
history requirements, and drug exclusions in dozens of criteria per 
trial. Manual screening at this scale is labor-intensive and 
error-prone.

Since 2023, the release of GPT-4 and the proliferation of open-source 
LLMs have prompted a wave of LLM-based matching systems, of which 
TrialGPT~\cite{jin2024trialgpt} in \emph{Nature Communications} is the 
most influential example. These systems report high accuracy on 
per-criterion eligibility classification and are beginning to be 
deployed in real recruitment workflows. A scoping review published in 
2025~\cite{chen2025scoping} identified twenty-four such systems 
published in 2024 alone. Yet the same review noted a striking common 
feature: none of these systems reports calibration metrics. No Expected 
Calibration Error, no Brier score, no reliability diagram, no 
calibration-in-the-large. The probabilistic reliability of LLM-based 
trial matching thus remains empirically undocumented at a moment when 
the technology is transitioning from research prototype to clinical 
deployment.

% ----------------------------------------------------------------------------
\section{Problem Statement}
\label{sec:intro-problem}

The absence of calibration reporting in LLM-based clinical trial 
matching is not a minor omission. Van Calster et 
al.~\cite{vancalster2019calibration}, in their influential BMC Medicine 
article, describe calibration as ``the Achilles heel of predictive 
analytics'' and argue that calibration is a prerequisite for the 
clinical deployment of any predictive model. The 2024 update of the 
TRIPOD guidelines, TRIPOD+AI~\cite{collins2024tripod}, formalizes this 
requirement: any AI-based clinical prediction model must report 
calibration curves, calibration-in-the-large, and calibration slopes. 
Failing to do so is grounds for methodological rejection by several 
leading medical journals.

The problem addressed in this thesis is therefore the following: 
LLM-based patient--trial matching systems are being deployed in 
clinical settings without any systematic audit of the probabilistic 
reliability of their outputs. This deployment carries three specific 
risks. First, overconfident predictions on wrong labels may direct 
patients toward inappropriate trials, or exclude patients from trials 
in which they would have been eligible. Second, the presence of 
systematic biases in evasion labels --- when a model reports 
``not enough information'' rather than committing to a decision --- may 
be undetectable through the standard accuracy metrics used to evaluate 
these systems. Third, the effect of common LLM post-training 
procedures, in particular reinforcement learning from human 
feedback~\cite{ouyang2022training}, on the calibration of medical 
predictions remains empirically undocumented. The present thesis 
addresses these risks through a systematic audit of twelve 
open-source LLM configurations.

% ----------------------------------------------------------------------------
\section{Research Question and Hypotheses}
\label{sec:intro-question}

The main research question addressed in this work is stated as follows:

\begin{quote}
\itshape
Are the eligibility scores produced by large language models in 
patient--clinical trial matching probabilistically reliable, and does 
reinforcement learning from human feedback affect this reliability?
\end{quote}

This question is decomposed into four falsifiable hypotheses, each 
associated with a numerical rejection threshold defined prior to the 
experimental phase:

\begin{itemize}
    \item \textbf{H1} predicts that current LLM pipelines are systematically 
    miscalibrated, with Expected Calibration Error exceeding the clinical 
    threshold of $0.05$.
    \item \textbf{H2} predicts that reinforcement learning from human 
    feedback degrades calibration, extending a finding of Kadavath et 
    al.~\cite{kadavath2022know} from general question answering to the 
    clinical domain. This is the central pre-registered contribution of 
    the thesis.
    \item \textbf{H3} predicts that model scores respect the ordinal 
    structure of the eligibility labels, with a concordance index above 
    $0.70$.
    \item \textbf{H4} predicts that confidence-based selective 
    classification improves macro-F1 by at least $0.02$ when coverage is 
    restricted to the top $80\%$ most confident cases.
\end{itemize}

The pre-registration of these hypotheses, together with the fixed 
rejection thresholds, aligns the study with the reporting standards of 
TRIPOD+AI~\cite{collins2024tripod} and protects the analysis against 
post-hoc adjustment of the criteria to fit the observed results.

% ----------------------------------------------------------------------------
\section{Contributions}
\label{sec:intro-contributions}

The thesis contributes to the calibration and clinical AI communities 
along four complementary directions.

The first contribution is methodological: the study provides a 
reusable protocol for auditing the calibration of LLM-based clinical 
prediction pipelines. The protocol combines a pre-registered 
hypothesis-testing framework, a log-probability scoring pipeline that 
produces well-defined probability distributions, a double-prompt 
design that treats prompt engineering as an experimental variable, and 
a composite utility score that integrates calibration, discrimination, 
and prediction diversity. Every design element is documented in 
Chapter~\ref{chap:methodology} and archived for independent 
reproduction (Appendix~\ref{app:code}).

The second contribution is confirmatory: the study tests the 
pre-registered hypothesis that RLHF degrades calibration on two 
independent model families (Mistral-7B and Llama-3.1-8B) and two 
independent prompt designs. All four permutation tests reject the null 
hypothesis at Bonferroni-corrected significance ($p < 0.003$), 
extending to the clinical domain a finding previously documented only 
in general question answering.

The third contribution is exploratory: the study documents twelve 
original findings that emerged from the deeper analysis. Three of 
these deserve particular emphasis. The enriched-prompt paradox (F3) 
shows that expert prompt engineering degrades calibration in five of 
six models. Mode collapse (F7 and F8) affects five of the twelve 
configurations, which concentrate more than $97\%$ of their 
predictions on a single label, and this collapse persists under 
constrained decoding. Complete failure on truly discriminative cases 
(F9) reveals that the two medical LLMs tested achieve accuracy at or 
below $0.4\%$ on the $238$ pairs requiring a genuine clinical decision 
--- despite their apparent superiority on global accuracy. These 
three findings jointly demonstrate that global accuracy and ECE, 
taken in isolation, are insufficient to characterize the clinical 
usability of an LLM classifier.

The fourth contribution is mechanistic: a deeper analysis of the 
log-probabilities reinterprets mode collapse as a property of the 
decision rule rather than a loss of knowledge. The study shows that 
the discriminative signal survives collapse (area under the curve up 
to $0.91$) but is masked by a response prior, identifies a scalar 
signal-to-prior ratio that predicts collapse across configurations 
(F11), and demonstrates that a simple re-thresholding rule recovers up 
to $+0.79$ macro-F1 where monotone recalibration provably cannot (F10). 
The same analysis separates recoverable collapse from an irrecoverable 
corruption of the signal, and both the collapse and its mechanism 
replicate on the independent TREC Clinical Trials corpus (F12). This 
contribution turns a descriptive failure mode into a diagnosable and 
often correctable property of the decoding step.

% ----------------------------------------------------------------------------
\section{Structure of the Thesis}
\label{sec:intro-structure}

The remainder of this thesis is organized as follows.

Chapter~\ref{chap:sota} reviews the state of the art at the 
intersection of clinical trial matching, probabilistic calibration, 
and medical large language models, and identifies the specific gap 
that motivates this work.

Chapter~\ref{chap:context} describes the academic framework of the 
project, the origin and evolution of the research question, the 
TrialGPT-Criterion-Annotations benchmark used as the empirical 
foundation, and the ethical and computational constraints that shape 
the study.

Chapter~\ref{chap:methodology} presents the experimental design in 
full detail: the four pre-registered hypotheses, the data preparation, 
the model and prompt selection, the log-probability scoring pipeline, 
the evaluation metrics, and the statistical analysis plan.

Chapter~\ref{chap:results} reports the empirical findings organized 
around the four pre-registered hypotheses and the twelve original 
findings. The chapter begins with the pre-LLM audit of the dataset, 
proceeds through the calibration audit and the RLHF effect, documents 
mode collapse and its reinterpretation as a decoding artefact, and 
closes with an external replication on the TREC Clinical Trials corpus.

Chapter~\ref{chap:conclusion} summarizes the outcomes, answers the 
research question, discusses the limitations of the study, and 
proposes future directions, including broader external validation 
beyond the TREC replication reported here, multi-agent adversarial 
debate as a mitigation for mode collapse, and mechanistic 
investigation of the origin of the response prior in medical LLMs.